\section{Related literature}\label{sec:literature}

This paper sits at the intersection of three strands of research. The first concerns the run dynamics and price stability of fiat-backed stablecoins; the second concerns the design, adoption, and cost of retail payment infrastructure; the third treats stablecoins specifically as payment instruments rather than as a financial-stability object. We connect the first two by pricing the run-dynamics object as a structural expected loss and carrying it into a corridor-level payments comparison, where it enters as a cost input rather than a systemic-risk concern. A standing review of the empirical literature on stablecoins is in \citet{ante_fiedler_etal_2023} and a ten-year retrospective in \citet{dionysopoulos_urquhart_2024}.

\paragraph{Stablecoin run dynamics and price stability.} The empirical antecedents closest to the calibration exercise here are \citet{anadu_etal_2023_runs}, who document flight-to-safety patterns across stablecoin types in the March 2023 episode and identify a break-the-buck threshold near par, and \citet{ma_zeng_zhang_2025}, who show that the concentration of arbitrage at the major issuers shapes the trade-off between secondary-market price stability and run risk. The role of public information about reserve quality on the equilibrium of the holders' coordination game is formalised in \citet{ahmed_aldasoro_duley_2025}, whose model is closely related to the global-game specification we calibrate. The money-view perspective on stablecoin par defence, framed around forward-market price absorption, is developed in \citet{aldasoro_mehrling_neilson_2023}. The theoretical complement to the global-game framework is \citet{li_mancini_griffoli_etal_2026}, who use a Diamond-Dybvig setup to characterise the social planner's optimal stablecoin backing rule and identify central bank reserve remuneration as a policy lever for run-risk reduction. Earlier theoretical work on the determinants of stablecoin stability includes \citet{routledge_zetlin_jones_2022}, \citet{liao_2021}, \citet{lyons_viswanath_natraj_2023}, and \citet{li_mayer_2024}. The empirical microstructure of stablecoin primary and secondary markets is documented in \citet{watsky_etal_2024} and \citet{du_sonawane_watsky_2025}. Methodologically, the model is a bank-run global game in the tradition of \citet{morris_shin_1998} and \citet{goldstein_pauzner_2005_demand_deposits}, adapted from the deposit contract to stablecoin redemption; coordination games are by now an established way to model stablecoin stability rather than new ground, so the contribution is in the use of the equilibrium, not the apparatus. A complementary global-game result locates stablecoin fragility on the settlement rail rather than the issuer's balance sheet: \citet{klee_etal_2026_fragility} show that the interaction of network externalities and blockchain congestion can generate runs even when reserves are perfectly safe. Our coordination game is on the same methodological footing but on a different axis: the latent state is the issuer fundamental rather than network congestion, and the object priced is the holder's expected loss rather than the run probability.

\paragraph{Payment systems innovation and cross-border policy.} The design and adoption of upgraded public domestic instant payment systems is surveyed in \citet{bech_shimizu_wong_2017}, \citet{boar_etal_2020}, and \citet{frost_etal_2024_fast_payments}; the case of PIX in Brazil is analysed as a paradigmatic instance of central-bank-led infrastructure provision in \citet{duarte_etal_2022_pix}. Cross-border retail payments and the friction stack associated with correspondent banking are documented in \citet{cpmi_2018_cross_border_retail}, with subsequent monitoring of progress in \citet{cpmi_brief10_2024}. The G20 cross-border payments policy framework is laid out in \citet{fsb_2020_roadmap}, refined in \citet{fsb_2023_priority_actions}, and updated through the most recent consolidated progress report in \citet{fsb_2025_progress}. The methodological foundation of the per-corridor remittance prices panel used in the empirical exercise is \citet{beck_martinez_peria_2011}. The economics of the public digital-money alternative against which private stablecoins are measured, including the upgraded public rails and central bank digital currency, is surveyed in \citet{ahnert_etal_2024_economics_cbdc}.

\paragraph{Stablecoins as payment instruments.} A growing literature treats stablecoins specifically through a payments lens rather than a financial-stability lens. The closest companion to the present paper, methodologically distinct but identically framed, is \citet{copestake_etal_2026}, who use high-frequency stock-price variation around the US GENIUS Act to estimate the implicit market expectation of stablecoin penetration in payments: they report an eighteen percent reduction in the market value of listed incumbent payment firms, larger for those focused on cross-border payments. \citet{ahmed_clouse_etal_2025} discuss the payment-system efficiency and financial-stability implications of the post-GENIUS regulatory framework. The international monetary and financial system implications of stablecoin adoption are surveyed in \citet{aldasoro_frost_ito_2026}, using Visa and Allium data to characterise real-economy stablecoin transaction volumes; the policy debate is framed in \citet{hernandez_de_cos_2026}. The relationship between stablecoins and money market funds and the implications for monetary policy are examined in \citet{bis_wp_1219}. The taxonomic question of what counts as a fiat-backed stablecoin is treated in \citet{bis_papers_141}, and \citet{adachi_etal_2020} provide an early regulatory and financial stability perspective for the global stablecoin segment.

The framework adopted here differs from \citet{copestake_etal_2026} in two respects: where they recover the implicit market expectation of stablecoin penetration from the cross-section of incumbent stock prices around a regulatory event, we recover the structural expected loss that a holder bears for using the stablecoin as a payment vehicle and apply it to a corridor-level cost comparison. The two methodologies are complementary: the implicit market expectation from one and the structural cost decomposition from the other speak to the same substantive question (will stablecoins capture a meaningful share of retail payments) from the equity-market and the user-cost sides respectively. The contribution here is also orthogonal to the systemic-risk strand of the literature: a regulator concerned with the macro-financial implications of stablecoin growth will draw on \citet{aldasoro_frost_ito_2026}, \citet{bis_wp_1219}, and \citet{anadu_etal_2023_runs}, while a regulator concerned with the choice and pricing of payment infrastructure will draw on the comparison developed in the present paper.
